\chapter{Submission below Metal}\label{ch:submission-below-metal}

Without Metal, a process reaches the GPU through GPU and shared-memory allocations, a set of launch words, a submission
record, a call into the kernel and a completion notification. The resident runtime that assembles these into a model
execution, and validates its results, is in \cref{ch:below-metal-native}.

A Metal application asks Metal to build a pipeline and encode a command buffer; Metal's driver library
(AGXMetal) then writes command pages and calls Apple's private \code{IOGPU} framework, which enters the kernel driver. A
below-Metal process does the same work itself: it allocates GPU memory and shared memory through the AGX driver's user
client, writes the code, the launch state and the command pages, builds the submission record, and calls the same \code{IOGPU}
submission function. The kernel driver, the GPU firmware and the hardware are unchanged.

The project's process owns the program bytes, the launch words, the command pages and the record.
IOGPU took ownership of the submitted callback blocks on every observed successful Submit; ownership on failure, timeout
or teardown is not established (\cref{sec:completion-callbacks-their}). The kernel driver and firmware own
everything after the trap.

Launches are known to work in the configurations whose layout was measured. Within the
measured graphs the words a launch needs are known by role, except for some known only by value, such as the
shader-state bits and the kernel-page word at offset 0x3c0; those are copied as captured. The runtime refuses every
configuration outside the envelope. Some of those might work, but none has been measured.

Interface terms:

\begin{itemize}
\item \emph{User client}: the IOKit connection to the AGX kernel driver. Its numbered entry points are \emph{selectors}.
\item \emph{Allocation n}: the n-th GPU memory allocation in a measured resource graph, created through selector 9. Each has a
  GPU virtual address (VA) and a CPU mapping at the same address.
\item \emph{Command pages}: two 16 KiB shared-memory regions per submission, created through selector 14: the \emph{kernel page}
  (kernel-command shared memory) and the \emph{segment page} (the segment list). Each has a shared-memory ID.
\item \emph{Launch state}: the per-launch words in measured allocations that select a program, its grid and its buffers.
\item \emph{Submit record}: the record handed to \code{IOGPUCommandQueueSubmitCommandBuffers}.
\item \emph{Measured graph}: a complete set of allocations, pages and launch words whose layout was captured from Metal once,
  then rebuilt from fields. Four sizes are in use: 3 MiB, 4 MiB, 128 MiB and 1 GiB.
\end{itemize}

Standing in this chapter: \emph{measured} means executed below Metal with a checked output and a control; \emph{observed} means
traced in a live process (Frida) or read from Apple's binary without executing a GPU command; \emph{static} means resolved by
reading the installed \code{IOGPU} code (\code{tools/g17submitinspect.c}).

\section{Apple-owned layers}\label{sec:stays-apple-s}

\begin{xltabular}{\linewidth}{@{}>{\hsize=0.59\hsize}X>{\hsize=1.26\hsize}X>{\hsize=1.15\hsize}X@{}}
\caption{What the below-Metal process writes and what remains Apple's, by layer.}\label{tab:apple-layers}\\
\toprule
Layer & Below Metal & Still Apple's \\
\midrule
\endfirsthead
\tablecontinued{3}
\toprule
Layer & Below Metal & Still Apple's \\
\midrule
\endhead
program & every instruction byte, from the project's compiler & encoding rules recovered with Apple's decoder \\*
resources and launch & allocations, command pages, launch words, bindings, the Submit record,
written by the project's host code & the formats themselves, recovered by measurement \\*
submission & the call into \code{IOGPUCommandQueueSubmitCommandBuffers} & \code{IOGPU.framework}, IOKit, the AGX kernel driver \\*
execution &  & the GPU firmware and the hardware \\*
\bottomrule
\end{xltabular}

The process never creates an \code{MTLDevice} or a Metal command buffer, and it checks that AGXMetal is not loaded
(MM~\mm{25.210.9}).

\section{Device, context and queue setup}\label{sec:opening-device-contexts}

\begin{itemize}
\item The AGX user client opens with no entitlement. A process with no Metal driver loaded creates the device and a
  context (measured; \code{tools/g17dispatch.c}, README rows 1–2).
\item Creating a command queue calls selectors 7, 16 and 28, in \code{IOGPUCommandQueueCreate} (observed in a Metal process that
  created only a device and a queue). Selector 28 registers the queue's completion notification queue (static). Calling
  the group a second time by hand ends in a selector-28 error (\code{docs/g17-first-dispatch-trials.md}, "Submission transport
  and record, resolved statically" and the preflight that follows it).
\item The sampler-source query (selector 261) returns a privilege error; nothing below depends on it (MM~\mm{25.143}).
\end{itemize}

\emph{Correction:} an earlier reading named selector 7 the submit and fence, and selector 14 a telemetry ring (MM~\mm{25.143}). The
static resolution places 7, 16 and 28 in queue creation and identifies selector 14 as the shared-memory allocator
(\cref{sec:memory-allocations-shared}), superseding the earlier attribution.

\section{Allocations and address windows}\label{sec:memory-allocations-shared}

\begin{itemize}
\item GPU allocations come from selector 9. Each request names an allocation class (a request \emph{type}), a size, and an
  address; the driver returns a buffer mapped at the same CPU and GPU virtual address (measured; MM~\mm{0.15}, README row 2).
\item The address window depends on the request type. The exact type value 0x58000000 places an allocation in the window
  at 0x6f\_00000000; other tested values do not reach the 0x1ea\_ window that Metal's command buffer names. Whether that
  window is the shader window is not shown (measured; MM~\mm{25.143.O}).
\item Shared memory comes from selector 14 (\code{IOGPUDeviceCreateDeviceShmem}). Its 16-byte output is
  \code{\{cpu\_address: u64, size: u32, shmem\_id: u32\}}. A submission uses two 16 KiB regions, the kernel page and the segment
  page, with IDs such as 2 and 1 (static and observed; trials, "Submission transport and record").
\item Allocations of 32 KiB, 8 MiB, 96 MiB and 1 GiB were created and CPU-mapped with sentinels and no Submit.
  This shows allocation and mapping only; GPU visibility was not established, so the runtime reuses measured graphs
  instead of a size formula (MM~\mm{25.210.3}).
\end{itemize}

\section{Code placement and the program counter}\label{sec:code-placement-program}

\begin{itemize}
\item Every program is uploaded into allocation 0, at GPU address 0x10000000000, and the allocation is treated as immutable
  after upload. Programs start at offset 0x6c0 and are 64-byte aligned (measured; MM~\mm{25.210.6}).
\item A launch's program counter is allocation 0's base plus the program's offset. It is written in pieces into allocation 23:
  the low 16 bits OR 7 at offset 0x40, bits 16–31 at offset 0x46, bits 32–47 at offset 0x48 (MM~\mm{25.210.6}).
\item The shader-state word at offset 0x42 of allocation 23 is written separately. A single 32-bit program-counter write once
  corrupted it (measured; MM~\mm{25.210.6}).
\item \emph{Not known:} the low program-counter tag bits outside the measured entry class.
\end{itemize}

\section{Launch state and command pages}\label{sec:launch-state-command}

\Cref{tab:launch-words} lists the words a launch needs, all written by the project's runner (MM~\mm{25.210.6}, unless noted).

\begin{xltabular}{\linewidth}{@{}>{\hsize=1.14\hsize}X>{\hsize=0.67\hsize}X>{\hsize=1.38\hsize}X>{\hsize=0.81\hsize}X@{}}
\caption{Launch words by location, with width, meaning and standing.}\label{tab:launch-words}\\
\toprule
Where & Width & Meaning & Standing \\
\midrule
\endfirsthead
\tablecontinued{4}
\toprule
Where & Width & Meaning & Standing \\
\midrule
\endhead
allocation 23, offsets 0x40, 0x46, 0x48 & 16 bits each & program counter: low half OR 7, bits 16–31, bits 32–47 & measured \\
allocation 23, offset 0x42 & 16 bits & shader state & measured value; bits unnamed \\
allocation 23, offset 0x38 & two 32-bit words & scratch/config word, then zero & measured value; bits unnamed \\
allocation 22 offset 0xa8; allocation 25 offset 0x10 & three 32-bit words each & grid x, y, z, written twice & measured \\
allocation 25, offset 0x1c & 32 bits & threads per threadgroup (32 in the measured classes) & measured \\
allocation 28, offset 0x1ba0 & four 64-bit pointers & program buffer pointers; at most three are admitted & measured \\
allocation 23, offset 0x06 & 16 bits & scaled metadata address coordinate: \code{(VA28 >> 13) \& 0xffff} & measured (see below) \\
allocation 25, offset 0x09 & 16 bits & scaled address coordinate: \code{(VA23 >> 14) \& 0xffff} & measured equality \\
kernel page, offset 0x3c0 & word & needed for the tensor contribution: clearing it removes the MMA result
while Submit and callbacks succeed & measured; meaning unknown \\
kernel page, offsets 0x3ec and 0x3f8 & words & kept as measured & unknown \\
\bottomrule
\end{xltabular}

Launch state comes in two classes. Scalar and packing programs use the state word 0x0e40 with the scratch word 0x0c000007. Tensor
projections and the cooperative LayerNorm use 0x0e58 with 0x0c00100f. Their individual bits are not named (MM~\mm{25.210.6}).

MM~\mm{25.210.6} lists the halfwords at allocation 23 offset 0x06 and allocation 25 offset 0x09
as "graph-identity correlates; no formula". The trial log fits both exactly at six graph sizes from 17 to 34 MiB, and in
the 4 MiB graph, as the shifted addresses above (\code{docs/g17-first-dispatch-trials.md}, the address-coordinate table). For
allocation 23's halfword the evidence is causal in one graph: changing only that halfword between the high and low
coordinate decided pass or fail. Treat the formula as established for the measured graphs; do not assume which firmware
step reads it, or that every allocation with that encoding has the same role.

In the 1 GiB class the complete 32 KiB shader record is copied into allocation 17 at offset 0x10000 and
selected by allocation 25's halfword (0x001a) at offset 0x09; the originally wrapped selector value (0x0048)
page-faulted and must not be reused. The 48 KiB binding metadata is copied from allocation 28 into allocation 17 and
selected by allocation 23's halfword (0x002c) at offset 0x06. The executor rewrites the selected low copy of the shader
record for each stage; writing the original high record after the copy does not change the program that runs
(measured; MM~\mm{25.210.3}).

\section{Submit record and IOGPU call}\label{sec:submit-record-call}

The Submit record's base layout is 48 bytes; Metal's observed records are 64 bytes, one per command buffer
(\cref{tab:submit-record}; static and observed; trials, "Submission transport and record").

\begin{xltabular}{\linewidth}{@{}lX@{}}
\caption{Fields of the Submit record by byte offset.}\label{tab:submit-record}\\
\toprule
Offset & Field \\
\midrule
\endfirsthead
\tablecontinued{2}
\toprule
Offset & Field \\
\midrule
\endhead
0x00, u32 & shared-memory ID of the kernel page \\*
0x04, u32 & shared-memory ID of the segment page \\*
0x08, u32 & optional sideband shared-memory ID \\*
0x10, u64 & scheduled callback block \\*
0x18, u64 & completed callback block \\*
0x20, u32 & optional debug shared-memory ID \\*
0x28, u64 & declared by the framework's metadata; meaning not
established \\*
0x30 on & possible driver extension; kept uninterpreted \\*
\bottomrule
\end{xltabular}

The runner must also meet these conditions (measured; MM~\mm{25.210.6}):

\begin{itemize}
\item \code{IOGPUDeviceGetNextGlobalTraceID} must return consecutive IDs for the trace-ID words in the command pages.
\item The coordination word at offset 0x24 of shared region 0 must hold the fire value 0x800000f0. The staged value 0xf0 was
  rejected.
\item Exactly one scheduled and one completed callback are required, each a fresh \code{Block\_copy}. Reusing blocks crashed the
  CPU notification dispatcher; the crash is in host code and says nothing about the GPU.
\end{itemize}

The call reaches the kernel in one of two ways (static; trials, "Submission transport and record"):

\begin{itemize}
\item A one-record Submit calls \code{IOConnectTrap4}: a kernel trap with trap index 0, the queue ID, the record size, the record
  pointer and an output pointer. The Metal reference runs observed exactly this, with count 1 and stride 64.
\item A multi-record Submit calls \code{IOConnectCallMethod} with selector 0x1d, passing the contiguous record array.
\end{itemize}

\emph{Correction:} MM~\mm{25.143} and README row 6a describe a Submit with no per-dispatch kernel call (shared-memory writes
only, a candidate store at Submit+0x11c). The static resolution shows each Submit enters the kernel through the trap
or the method call above; the store at Submit+0x11c is an output-argument write (\cref{app:a}).

\emph{Open: what consumes the fire value.} In a Metal process, a control that never set the ready bit still executed
(MM~\mm{25.143}). In the pure process, the staged value 0xf0 was rejected and only 0x800000f0 works (MM~\mm{25.210.6}). The two
observations come from different paths and do not settle what consumes the word. \emph{Safe assumption:} a pure process
writes 0x800000f0 before Submit. \emph{Not safe:} treating that write as the mechanism that starts the GPU.

\section{Completion and callback ownership}\label{sec:completion-callbacks-their}

\begin{itemize}
\item Completion arrives as a notification. \code{IOGPUNotificationQueueDispatchAvailableCompletionNotifications} dequeues
  40-byte notifications (\code{IODataQueueDequeue}) and invokes the referenced callback block. In the Metal reference runs,
  two notifications per submission carried, as their first 64-bit values, exactly the record's scheduled and completed
  callback pointers, in that order (observed; trials, "Submission transport and record").
\item In the MiniLM FFN block, a warm request spends 2.385~ms from Submit to the completed-callback entry and 0.0495~ms
  from callback entry to the waiter's return (measured; MM~\mm{25.210.5}). The first interval mixes GPU
  execution, driver work and completion delivery and is not GPU time.
\item Callback ownership was observed on macOS 26.6.2 (MM~\mm{25.210.9a}). On the successful count-one path, each of six observed callback
  blocks was invoked, returned, released by the notification dispatcher, and freed through \code{\_Block\_release}. No caller
  release occurred, and no block accumulated. \emph{Policy this supports:} allocate fresh copied blocks per Submit, hand them
  to IOGPU, and neither release nor reuse them. \emph{Not known:} ownership on failure, timeout, or queue teardown with
  pending callbacks, and behaviour on other OS versions.
\end{itemize}

\section{Multi-record Submits and dependent dispatches}\label{sec:several-records-per}

\begin{itemize}
\item Several records in one Submit execute. Count-two and count-three Submits ran distinct programs and bindings into
  separate outputs, and later trials reached seven separately bound records, the seventh predicted without a Metal
  capture (measured; trials, "Two command records in one pure Submit" through "Seventh record predicted").
\item Dependent stages, as the resident runtime runs them, are consecutive count-one Submits: 7 for the MiniLM FFN, 916
  for one Qwen decode token (measured; MM~\mm{25.210.3}, MM~\mm{25.210.5}). A two-stage dependent pipeline, with and without a
  program switch, passed in the trials (measured; trials, "Dependent two-stage pure GPU pipeline").
\item \emph{Not built:} many dependent dispatches in one Submit with explicit barriers between them. Metal's serial-barrier word
  and fence semantics are decoded (README row 10), but the runtime does not use them.
\end{itemize}

\section{Faults, hangs and recovery}\label{sec:faults-hangs-recovery}

\begin{itemize}
\item Malformed kernels have hung the GPU on this machine. Killing the host process does not cancel a submitted command
  buffer; recovery from that hang was a reboot (rule 35 of MM~\mm{17}).
\item The GPU's recovery counter and event reports are observable before and after a run. The resident runtime refuses a run
  if either changes (MM~\mm{25.210.9}).
\item The runtime bounds the wait for each Submit's callbacks to five seconds (MM~\mm{25.210.6}).
\item \emph{Not available:} detecting or recovering from a hang beyond that bound, and resetting the GPU (README row 11).
\end{itemize}

\section{Scope of the measured graphs}\label{sec:measured-graphs-generalizes}

\begin{xltabular}{\linewidth}{@{}l>{\hsize=1.27\hsize}X>{\hsize=0.73\hsize}X@{}}
\caption{Measured resource graphs by size, with the work each carries and its source.}\label{tab:measured-graphs}\\
\toprule
Graph & Used for & Source \\
\midrule
\endfirsthead
\tablecontinued{3}
\toprule
Graph & Used for & Source \\
\midrule
\endhead
3 MiB & the MiniLM FFN and attention blocks & MM~\mm{25.210.4}, MM~\mm{25.210.6} \\*
4 MiB & the connected layer (allocation 20 enlarged to 4 MiB, allocation 21
moved) & MM~\mm{25.210.6} \\*
128 MiB & the MiniLM encoder & MM~\mm{25.210.3} \\*
1 GiB & the Qwen2.5-0.5B decoder & MM~\mm{25.210.3} \\*
\bottomrule
\end{xltabular}

Outside these graphs the program-counter placement, the grid and threadgroup words, up to three buffer pointers, the
two state classes, the record layout, and the address-coordinate halfwords within the tested range still apply. Not
yet carried over:

\begin{itemize}
\item arbitrary graph allocation and publication (the runtime recovers only the request and page sets of the measured
  graphs, and refuses every other size);
\item a fourth buffer pointer (earlier controls wrote 0 of 32 outputs);
\item new shapes, data types, simdgroup counts, OS versions or GPU generations;
\item GPU-only timestamps.
\end{itemize}

\textbf{Evidence.} MM~\mm{0.15}, MM~\mm{25.143} (with 25.143.N and 25.143.O), MM~\mm{25.147}, MM~\mm{25.150}, MM~\mm{25.210} (25.210.3, 25.210.5,
25.210.6, 25.210.9, 25.210.9a); \code{docs/g17-first-dispatch-trials.md}; README rows 1–16.
